\documentclass[10pt,a4paper]{amsart}
\usepackage[margin=19mm]{geometry}
\usepackage{amsmath,amssymb,mathtools}
\usepackage[scr=rsfs,cal=euler]{mathalfa}
\usepackage{xfrac}
\usepackage[unicode,hidelinks,bookmarks=false]{hyperref}
\newcommand{\ie}{i.e.\ }
\newcommand{\cf}{cf.\ }
\newcommand{\resp}{respectively\ }
\newcommand{\Subsection}{\subsection}
\providecommand{\nobreakdash}{\nobreak\hbox{-}}
\setlength{\emergencystretch}{2em}
\multlinegap0pt
\usepackage{bm}
\usepackage{bbm}
\usepackage{dsfont}
\usepackage{xspace}
\usepackage{cancel}
\usepackage{mathtools}
\usepackage{amsthm}
\makeatletter
\@ifundefined{definition}{\newtheorem{definition}{Definition}}{}
\@ifundefined{theorem}{\newtheorem{theorem}{Theorem}}{}
\makeatother
\title[On Spectral Invariant Dense Subalgebras of Uniform Roe Algeb]{On Spectral Invariant Dense Subalgebras of Uniform Roe Algebras with Subexponential Growth}
\author{Siqi Jiang, Xianjin Wang}
\date{Source: arXiv:2507.17322v1 (CC BY 4.0)}
\begin{document}
\maketitle

\noindent\emph{Source and licence.} On Spectral Invariant Dense Subalgebras of Uniform Roe Algebras with Subexponential Growth. Version arXiv:2507.17322v1 (CC BY 4.0), distributed under the Creative Commons Attribution 4.0 International licence, as recorded in the arXiv licence metadata for this exact version and shown on the abstract page https://arxiv.org/abs/2507.17322v1. Full rights notice: licenses/arxiv-2507.17322-CC-BY-4.0.txt.\par\smallskip

\noindent\textit{Editorial scope.} Counted unit: the Theorem whose source label is \texttt{thm4.2} in the article (source0.tex lines 250--252, source label \texttt{thm4.2}), delivered together with the article's own complete original proof of it (source0.tex lines 254--395, one \texttt{proof} environment of 142 lines). No mathematics is written, repaired or re-derived: statement and proof are the article's own text line for line. Required context retained uncounted: w1 (lines 171--195); w22 (lines 177--179); w2 (lines 181--184). Every definition, display and label of those lines is retained unchanged. Omitted from this package and not redistributed: 1--170, 180--180, 185--249, 253--253, 396--703. Nothing inside a retained excerpt is omitted or altered: every retained body is delivered exactly as the article's lines are written, so no omission receipt of any kind accompanies this package. Citations: the retained excerpts carry no citation command at all, so no bibliography entry of the article is transcribed and no reference is redistributed. Reference adaptation: none was needed and none was made; every reference inside the delivered unit resolves inside it, so no unresolved reference or citation is produced. Printed slips kept literal: no printed word, symbol, hypothesis or number of the counted statement or of its proof was altered. Layout and class: the article's own class is \texttt{elsarticle} with packages including fontenc, hyperref, mathptmx, amsmath, amssymb, amsfonts, amsthm, enumerate, color, mathbbol, bm, tensor, booktabs, cases; the wrapper is the lane's pinned \texttt{amsart} editorial wrapper at 10pt on A4, which restates the theorem-like environments the retained text needs and adds the article's own macro definitions that the retained text needs (none), with extra wrapper packages (bm, bbm, dsfont, xspace, cancel, mathtools). The delivered numbering is the unit's own. Assets: no retained span contains a figure environment, a graphics-inclusion command, a PDF-page inclusion, a table or a TikZ picture; no image file is copied into this package, the package declares no asset dependency and no image file is redistributed with it. Licence: version arXiv:2507.17322v1 of this article is distributed under the Creative Commons Attribution 4.0 International licence, as recorded in the arXiv licence metadata for this exact version and shown on the abstract page https://arxiv.org/abs/2507.17322v1. A full rights notice is delivered at licenses/arxiv-2507.17322-CC-BY-4.0.txt. Characters: every retained excerpt is pure ASCII, so the delivered engine, XeLaTeX with its default Latin Modern text font, reports no missing character.
	\begin{definition}
		Let $1\le p,r\le \infty$. We say that a weight $\omega $ is $\left (p,r  \right ) $-{\it admissible} if there exist another weight $v$ and two positive constants $D\in \left ( 0,\infty  \right ) $ and
		$ \theta \in \left ( 0,1 \right ) $ such that
		\begin{equation}\label{w1}
			w\left ( x,y \right ) \le  D\left ( w\left ( x,z \right ) v\left ( z,y \right ) +v\left ( x,z \right ) w\left ( z,y \right )  \right )~\text{for~all}~ x,y,z\in G,
		\end{equation}
		\begin{equation}\label{w22}
			\sup_{x\in G}\left \| (vw^{-1})\left ( x,\cdot \right ) \right \|_{p^{_{'}}} +\sup_{y\in G}\left \|( vw^{-1})\left ( \cdot,y \right ) \right \| _{p^{_{'}}}\le D,
		\end{equation}
		and
		\begin{equation}\label{w2}
			\inf _{\tau >0}a_{r' }\left ( \tau  \right ) +b_{p' }\left ( \tau  \right )t                                                                                                                                                                                                                                                                                                                                                                                                                                                                                                                                                                                                                                                                                                                                                                                                                                                                                                                                                                                                                                                                                                                                                                                                                                                                                                                                                                                                                                                                                                                                                                                                                                                                                                                                                                                                                                                                                                                                                                                                                                                                                                                                                                                                                                                                                                                                                                                                                                                                                                                                                                                                                                                                                                                                                                                                                                                                                                                                                                                                                                                                                                                                                                                                                                                                                                                                                                                                                                                                                                                                                                                                                                                                                                                                                                                                                                                                                                                                                                                                                                                                                                                                                                                                                                                                                                                                                                                                                                                                                                                                                                                                                                                                                                                                                                                                                                                                                                                                                                                                                                                                                                                                                                                                                                                                                                                                                                                                                                                                                                                                                                                                                                                                                                                                                                                                                                                                                                                                                                                                                                                                                                                                                                                                                                                                                                                                                                                    
			\le Dt^\theta ~ \text{for~all}~ t\ge 1,
		\end{equation} 
		where ${p}' =p/\left(p-1 \right)$, $r' =r/\left ( r-1 \right )$,
		\begin{equation*}\label{w3}
			a_{r'}(\tau)=\sup_{x\in G}\left \| v\left ( x,\cdot \right )\chi _{B\left ( x,\tau  \right )\left ( \cdot  \right )  }  \right \|_{r^{'}} +\sup_{y\in G}\left \| v\left ( \cdot,y \right )\chi _{B\left (y, \tau  \right )\left ( \cdot  \right )  } \right \|_{r^{'}},
		\end{equation*}
		\begin{equation*}
			b_{p^{'}}(\tau)=\sup_{x\in G}\left \| (vw^{-1})\left ( x,\cdot \right )\chi _{X\setminus B\left ( x,\tau  \right )\left ( \cdot  \right )  }  \right \|_{p^{'}} +\sup_{y\in G}\left \|( vw^{-1})\left ( \cdot,y \right )\chi _{X\setminus B\left (y, \tau  \right )\left ( \cdot  \right )  } \right \|_{p^{'}},
		\end{equation*}
		$\chi _{E}$ is the characteristic function on the set $E$, and $\left \| \cdot  \right \| _{p}$ is the norm on $\ell^{p}$.
		
		Unless stated otherwise, in this paper, $p=2$.
	\end{definition}

	\begin{theorem}\label{thm4.2}
		Let $G$ be a countable discrete group with a proper length function $l$. If $G$ satisfies property $P$,  the subexponential weight $w(x,y)=\exp(\alpha \rho_{l }(x,y)^{\beta}$ is a $(2,r)$-admissible weight.
	\end{theorem}

	\begin{proof}
		
		In order to prove $w$ is an admissible weight satisfying (\ref{w1})-(\ref{w2}), the proof is carried out in several steps.
		
		\textbf{Step 1.}  The weights $w$ and $v$ satisfy (\ref{w1}), i.e.,
		$$w(x,y)\le D\left ( w(x,z)v(z,y)+ v(x,z)w(z,y)\right)$$ for all $x,y,z\in G$ with $D=1$.
		
		\noindent We recall the form of weight $w$ and define the weight $v$ as follows
		\begin{equation}
			\label{adw1}
			w(x,y)=\exp(\alpha \rho_{l }(x,y)^{\beta})\quad \text{for~all}~ x,y\in G,
		\end{equation}
		\begin{equation}\label{adv1}
			v(x,y)=\exp(\alpha (2^{\beta }-1) \rho_{l}(x,y)^{\beta}) \quad \text{for~all}~ x,y\in G,
		\end{equation}
		with $\alpha\in (0,\infty),\beta \in (0,1)$.
		Note that the following inequality holds
		$$ 1\leq s^\beta +(2^\beta -1)(1-s)^\beta\quad\text{for~all}~\tfrac{1}{2}\leq s \leq 1.$$
		Let 
		$$
		s=\begin{cases}
			\frac{\rho_l(x,z)}{\rho_l(x,z)+\rho_l(z,y)}, &\text{ if } \rho_l(x,z)\geq \rho_l(z,y)\\
			\frac{\rho_l(z,y)}{\rho_l(x,z)+\rho_l(z,y)}, &\text{ if } \rho_l(x,z)< \rho_l(z,y)
		\end{cases}.
		$$
		If $ \rho_l(x,z)\geq \rho_l(z,y)$, we have
		\begin{align*}
			w(x,y)&=\exp(\alpha \rho_{l }(x,y)^{\beta})\le \exp(\alpha( \rho_{l }(x,z)+\rho_{l }(z,y))^{\beta }) \\
			&\leq\exp(\alpha (s^\beta+(2^\beta-1)(1-s)^\beta)( \rho_{l}(x,z)+\rho_{l}(z,y))^{\beta})\\
			&=\exp(\alpha\rho_l(x,z)^\beta +(2^\beta-1)\rho_l(z,y)^\beta)=w(x,z)v(z,y).
		\end{align*}
		Similarly, if $\rho_l(x,z)< \rho_l(z,y)$, we have
		$$w(x,y)\leq v(x,z)w(z,y)\quad\text{for~all}~x,y,z\in G.$$
		Hence,
		$$w(x,y)\leq D\left(w(x,z)v(z,y)+v(x,z)w(z,y)\right)$$
		for all $x,y,z\in G$ with $D=1$.
		Thus, the weights $w$ and $v$ satisfy (\ref{w1}).
		
		\textbf{Step2.} The weights $w$ and $v$ satisfy (\ref{w22}), i.e.,
		\begin{align*}
			\sup_{x\in G}\left \| vw^{-1}(x,\cdot ) \right \| _{2}+\sup_{y\in G}\left \| vw^{-1}(\cdot,y ) \right \| _{2}<\infty.
		\end{align*}
		Firstly, we estimate $\sup_{x\in G}\left \| vw^{-1}(x,\cdot ) \right \| _{2}$.
		
		Based on {\it property P} of the group $G$, we know that for any $0<\alpha'<\alpha(2-2^{\beta})$ and $0<\beta'<\beta<1$, there exists $C_{\alpha' ,\beta' }$ such that
		\[ \left | B(x,\tau) \right | \le C_{\alpha' ,\beta' }\exp(\alpha' \tau^{\beta' })%=C\exp(\alpha' \tau^{\beta' })
		\quad\text{for~ all}~  x\in G~\text{and
		}~\tau\ge1.\]
		By the definition of the weights $v$ and $w$ , we have
		\begin{align}\label{vw}
			vw^{-1}(x,y)=\exp(-\alpha (2-2^{\beta }) \rho_{l}(x,y)^{\beta}).
		\end{align}
		Since 
		\begin{gather}\label{2.7'}
			\sup_{x\in G}\left \| vw^{-1}\left ( x,\cdot \right )\chi _{B\left ( x,\tau  \right )\left ( \cdot  \right )  }  \right \|_{2}=\sup_{x\in G}(  (	vw^{-1}(x,y)    )^{2} \left | B(x,\tau ) \right |)^\frac{1}{2},
		\end{gather}
		and
		\begin{gather}\label{2.7}
			\sup_{x\in G}\left \| vw^{-1}\left ( x,\cdot \right )\chi _{X\setminus B\left ( x,\tau  \right )\left ( \cdot  \right )  }  \right \|_{2}=(\sum_{j=0}^{\infty } \sum _{2^{j}\tau \le\rho_l (x,y)< 2^{j+1}\tau} ( vw^{-1}(x,y)  )^{2})^\frac{1}{2}.
		\end{gather}
		It is worth noting that the above results are also true for $y$.
		
		\noindent Combing  (\ref{vw}),(\ref{2.7'}) and (\ref{2.7}), we have
		\begin{align*}%\label{sub-2.51}
			&	\sup_{x\in G}\left \| vw^{-1}\left ( x,\cdot \right )\chi _{B\left ( x,\tau  \right )\left ( \cdot  \right )  }  \right \|_{2}= \sup_{x\in G}(\sum _{\rho_{l}(x,y)< \tau }\exp(-2\alpha (2-2^{\beta })\rho_{l }(x,y)^{\beta })^\frac{1}{2}\notag\\
			&\le C\sup_{x\in G}(\exp(-2\alpha (2-2^{\beta })\tau ^{\beta})\left| B(x,\tau)\right|)^\frac{1}{2}\notag\\
			&\le C\sup_{x\in G}(\exp(-2\alpha (2-2^{\beta })\tau ^{\beta}+\alpha'\tau^{\beta'}))^\frac{1}{2}\le C
		\end{align*}
		and
		\begin{align}\label{sub-2.52}
			&\sup_{x\in G}\left \| (vw^{-1})\left ( x,\cdot \right )\chi _{X\setminus B\left ( x,\tau  \right )\left ( \cdot  \right )  }  \right \|_{2}\leq\Big(\sum_{j=0}^{\infty} (\exp (-2\alpha (2-2^{\beta})(2^{j}\tau)^{\beta})) \left| B(x,2^{j+1}\tau)\right|\Big)^\frac{1}{2} \notag \\
			&\leq C\Big(\sum_{j=0}^{\infty } \exp(-2\alpha(2-2^{\beta}) 2^{j\beta}\tau^{\beta }+\alpha'2^{(j+1)\beta'}\tau ^{\beta'} )\Big)^\frac{1}{2}\notag \\
			&\leq C\Bigl(\sum_{j=0}^{\infty }\exp( -2(\alpha(2-2^{\beta}) -\alpha')2^{j\beta}\tau^{\beta} )\Bigr)^{\frac{1}{2}}\notag \\
			&=C \Bigl(\sum_{j=0}^{\infty }\exp(-\alpha(2-2^\beta)\tau^\beta[2(1-\tfrac{\alpha'}{\alpha(2-2^{\beta})})]2^{j\beta} \Bigr)^{\frac{1}{2}} \notag\\
			&=C \Bigl(\sum_{j=0}^\infty q^{2\bigl(1-\tfrac{\alpha'}{\alpha(2-2^{\beta})}\bigr)2^{j\beta}}\Bigr)^{\frac{1}{2}}\leq C  (\sum_{j=0}^\infty q^{2^{j+1}})^{\frac{1}{2}} \notag \\
			&\leq C (\sum_{j=1}^{\infty}q^{2j}  )^{\frac{1}{2}}\leq C\sqrt{\frac{q^2}{1-q^2}}\leq C'q ,
		\end{align}
		where $q=\exp(-\alpha(2-2^{\beta})\tau^\beta)<1$, and note $\tau\geq 1$, the $\frac{1}{\sqrt{1-q^2}}$ indeed has upper bound.
		
		Then, we obtain 
		\begin{align*}
			\sup_{x\in G}\left \| vw^{-1}(x,\cdot ) \right\| _{2}\le (\ref{2.7'})+(\ref{2.7})
			&\le C+C'\sum_{j=0}^{\infty } \exp(-\alpha(2-2^{\beta}) \tau^{\beta } )<\infty.
		\end{align*}
		%By the property $P$,we know for some $\beta'<\beta$ and $\alpha'<\alpha(2-2^{\beta})$, there exists $C_{\alpha',\beta'}$ such that
		%	\[ \left | B(x,r) \right | \le C_{\alpha' ,\beta' }\exp(\alpha' r^{\beta' })\quad\text{for~ all}~  x\in G~\text{and	}~r>0.\]
		%where $d(G,\rho_l)>0$, $\beta'<\beta$ and $\alpha'<\alpha(2-2^{\beta})$.
		Similarly, we  obtain 
		\begin{equation*}
			\sup_{y\in G}\left \| vw^{-1}(\cdot,y ) \right \| _{2}<\infty.
		\end{equation*}
		Thus, we have the weights $w$ and $v$ satisfy (\ref{w22}).
		%\[\sup_{x\in G}\left \| vw^{-1}(x,\cdot ) \right \| _{2}+\sup_{y\in G}\left \| vw^{-1}(\cdot,y ) \right \| _{2}<\infty,\]
		%which means the weights $w$ and $v$ satisfy (\ref{w22}).
		
		\textbf{Step 3.}  The weights $w$ and $v$ satisfy (\ref{w2}), i.e.,
		\begin{align*}
			\inf_{\tau\ge 1} a_{r'}(\tau ) +b_{p'}(\tau ) \cdot t\le Ct^{\theta}\quad \text{for~all} ~t\ge 1~\text{and}~\theta\in (\frac{1}{3-2^{\beta}},1).
		\end{align*}
		Firstly, we get the following estimate
		\begin{align}\label{sub2.5}
			\sup_{x\in G}\left \| v\left ( x,\cdot \right )\chi _{B\left ( x,\tau  \right )\left ( \cdot  \right )  }  \right \|_{r'} 
			%&=\sup_{x\in G}\Big[ \sum_{y\in G,\rho_{l }(x,y)\le \tau } \Big ( 	\exp(\alpha (2^{\beta }-1) \rho_{l }(x,y)^{\beta})\Big )^{r'} \Big]^\frac{1}{r'}\notag\\
			&\le\sup_{x\in G}( \exp(r'\alpha (2^{\beta }-1) \rho_{l }(x,y)^{\beta})\left| B(x,\tau)\right|) ^\frac{1}{r'}\notag\\
			%&\le C\left ( \exp	\left(\alpha (2^{\beta }-1)\tau^\beta \right )\cdot\exp(\alpha'\tau ^{\beta'})^\frac{1}{r'}\right )\notag\\
			&\le C\exp(\alpha (2^{\beta }-1)\tau^\beta+\alpha'\tau^{\beta'}/r')\notag\\
			&\le C\exp(\alpha (2^{\beta }-1)\tau^\beta+\alpha'\tau^{\beta'})\le C\exp(\alpha \tau^\beta),
		\end{align}
		where $r'\ge 1$. 
		
		\noindent Then, by applying the inequalities (\ref{sub-2.52}) and (\ref{sub2.5}), we derive that
		%$$a_{r'}(\tau )\le C\exp\left(\alpha \tau^\beta\right)~\text{and}~b_{p'}(\tau )\le C\exp(-\alpha(2-2^{\beta}) \tau^{\beta } ).$$
		%Hence, we get
		\begin{equation}
			\label{e:key_exp}
			\inf_{\tau\ge 1} a_{r'}(\tau ) +b_{p'}(\tau ) t
			\le C\inf_{\tau\ge 1} [\exp(\alpha \tau^\beta) + \exp(-\alpha(2-2^{\beta}) \tau^{\beta } )\cdot t].
		\end{equation}
		%Indeed, let $\tau=1+f(t)$ with $f(t)$ determined later, it is enough to prove
		We assert that there is a function $f(t)$ satisfies that		
		\begin{align}
			\exp(\alpha(1+f(t))^\beta)&\le Ct^{1/(3-2^\beta)}\label{e3.15}\\
			\exp(-\alpha(2-2^\beta)(1+f(t))^\beta)\cdot t&\le Ct^{1/(3-2^\beta)}\label{e3.16}.
		\end{align}
		Indeed, 
		\begin{align*}
			(\ref{e3.15}) &\iff \exp(\alpha(1+f(t))^\beta)\leq \exp(1/(3-2^\beta)\ln t+\ln C)\\
			&\iff \alpha(1+f(t))^\beta \leq \frac{\ln t}{3-2^\beta}+\ln C,
		\end{align*}
		and 
		\begin{align*}
			(\ref{e3.16})&\iff t^{\frac{2-2^\beta}{3-2^\beta}}\leq C\left[ \exp(\alpha(1+f(t))^\beta)\right]^{2-2^\beta}\\
			&\iff \exp\left(\frac{1}{3-2^\beta}\ln t\right)\leq C^{\frac{1}{2-2^\beta}}\exp(\alpha(1+f(t))^\beta)\\
			&\iff \frac{1}{3-2^\beta}\ln t\leq \alpha(1 + f(t))^\beta +\frac{1}{2-2^\beta}\ln C.
		\end{align*}
		Setting $\frac{1}{3-2^\beta}\ln t= \alpha(1 + f(t))^\beta$, i.e., 
		$$f(t)= \left(\frac{\ln t}{\alpha(3-2^\beta)}\right)^{\frac{1}{\beta}}-1,$$
		the desired inequalities (\ref{e3.15}) and (\ref{e3.16}) hold.
		Let $\tau=1+f(t)$ in (\ref{e:key_exp}), we get 
		\[\inf_{\tau\ge 1} a_{r'}(\tau ) +b_{p'}(\tau ) \cdot t\le Ct^{ 1/(3-2^{\beta })}\le Ct^{\theta}\quad \text{for~all}~t\ge 1~,\theta\in (\frac{1}{3-2^{\beta}},1),\]
		which means the weights $w$ and $v$ satisfy (\ref{w2}).
	\end{proof}

\end{document}
